\documentclass[11pt]{article}
\usepackage[review]{acl}

\usepackage{times}
\usepackage{latexsym}
\usepackage[T1]{fontenc}
\usepackage[utf8]{inputenc}
\usepackage{microtype}
\usepackage{inconsolata}
\usepackage{graphicx}
\usepackage{booktabs}
\usepackage{tabularx}
\usepackage{amsmath}
\usepackage{amssymb}
\usepackage{xcolor}
\usepackage{multirow}

% Shorthand macros
\newcommand{\splan}{$S_{\mathrm{plan}}$}
\newcommand{\splanmath}{S_{\mathrm{plan}}}
\newcommand{\sllm}{$S_{\mathrm{LLM}}$}
\newcommand{\Ra}{$R_a$}
\newcommand{\Ca}{$C_a$}
\newcommand{\Ramath}{R_a}
\newcommand{\Camath}{C_a}

\title{\textit{S\textsubscript{plan}}: Compact Learned Symbolic Planning Tokens\\
for Music-Aware Audio Language Models}

\author{
  Dabin Kim$^*$ \quad Minhee Lee$^*$ \quad Ji Jiaxian$^*$ \quad Aner Zheng$^*$ \\
  Graduate School of Cultural Technology, KAIST \\
  \texttt{\{dabinkim, minheelee, jiaxian, anerzheng\}@kaist.ac.kr}
}

\begin{document}
\maketitle

\begin{abstract}
We ask three questions about symbolic music representations in audio language models:
(1) does a compact learned symbolic stream reduce acoustic token uncertainty more than
a shuffled or dummy stream of the same size?
(2) is that benefit specific to content alignment, or merely to token capacity?
(3) does the symbolic stream encode music-structural attributes recoverable by probing?

We address these with \splan{}, a compact symbolic planning token that compresses
MIDI-derived musical structure into a fixed-length residual vector quantization (RVQ)
stream inserted between the reasoning and acoustic token streams of UniAudio 2.0.
Answering (1) and (2): across three independent experiments using aligned / shuffled /
dummy controls, aligned \splan{} consistently reduces $\Camath$ cross-entropy below
both corrupted controls---most rigorously in EXP034B (test CE $5.600$ vs.\ $5.614$
shuffled and dummy, $5.859$ reason-only), confirming an alignment-specific effect of
$\Delta = -0.014$ on top of a capacity effect of $\Delta = -0.245$.
Answering (3): MIR probing (EXP031A) recovers weak key accuracy $0.557$ ($+0.431$
over the zero-feature baseline) and pitch-class macro-F1 $0.854$ ($+0.024$), while
meter and instrument micro-F1 are prior-dominated and excluded from claims.
We further document two negative results (side-channel insertion and adapter gating)
that establish a stream-native interface as a necessary design constraint.
The central claim is conservative but supported: \splan{} is a compact symbolic codec
with acoustic-token utility and recoverable music-structural content under controlled
Slakh experiments.
\end{abstract}

%% ============================================================
\section{Introduction}
%% ============================================================

Unified audio language models (LALMs)~\citep{yang2024uniaudio,yang2026uniaudio2,rubenstein2023audiopalm,borsos2023audiolm}
predict discrete audio tokens autoregressively and have achieved strong performance
across diverse tasks.  Music, however, poses a persistent gap: the structural properties
most salient to musicians and listeners---harmonic progressions, rhythmic regularity,
instrument co-occurrence, phrase boundaries---are encoded in MIDI symbolic representations,
not waveform statistics.  While MIDI-to-audio synthesis pipelines use explicit symbolic
intermediates, integrating those representations into LALMs in a \emph{controlled,
verifiable} way has remained open.

\paragraph{What we are trying to prove.}
We organize our work around three research questions (RQs):

\begin{description}
  \item[RQ1 (Acoustic Utility):] Does a compact learned symbolic stream, aligned to
    an audio window's own MIDI, reduce acoustic token prediction uncertainty more than
    a shuffled or dummy stream of the same length?
  \item[RQ2 (Alignment Specificity):] Is the benefit from alignment-specific content,
    or from token capacity alone?  We test this by requiring aligned to beat
    \emph{both} shuffled (content broken, statistics preserved) and dummy (zero content,
    same length) controls.
  \item[RQ3 (Music-Structural Content):] Does the symbolic stream encode
    musically meaningful attributes recoverable by supervised probing?
    We test pitch, key, meter, density, and instrument recognition, each compared
    against a zero-feature baseline that exposes dataset priors.
\end{description}

Answering RQ1 affirmatively is not automatic: language models are good at ignoring
spurious context.  EXP027 demonstrates exactly this failure mode---aligned and shuffled
MIDI insertions produce nearly identical predictions when the interface does not match
the model's positional encoding.  Our positive result in RQ1 is thus contingent on the
correct interface design established through EXP027 and EXP028.

Answering RQ2 is important because any fixed-length stream adds context capacity.
A positive result on RQ1 alone is weak evidence if an equal-length dummy stream achieves
the same gain.  Our multi-condition design forces the claim: aligned must beat both
shuffled and dummy, or we cannot attribute the gain to symbolic content.

Answering RQ3 connects acoustic utility to an interpretable mechanism.  A symbolic stream
that reduces CE but encodes nothing auditable could be acting through any number of
latent channels; MIR probing makes the representational content explicit and falsifiable.

\paragraph{How we answer them.}
We introduce \splan{}, a \emph{learned symbolic planning token} that compresses
MIDI-derived features into a compact, fixed-length, multi-stream RVQ representation
and inserts it between the reasoning tokens $\Ramath$ and semantic acoustic tokens
$\Camath$ of UniAudio 2.0~\citep{yang2026uniaudio2}, forming
$\Ramath \to \splanmath \to \Camath$.
We evaluate through a systematic ladder of four experiments using aligned / shuffled /
dummy controls at each stage, plus MIR probing with a zero-feature baseline.

\paragraph{Contributions.}
\begin{enumerate}
  \item A low-rate RVQ symbolic codec (\splan{}) that compresses MIDI-derived frame
    features into $4 \times 30$ tokens per 30-second window---a $10\times$--$50\times$
    reduction versus raw MIDI tokenizations.
  \item \textbf{A positive answer to RQ1:} aligned \splan{} reduces $\Camath$
    cross-entropy consistently across EXP025, EXP029, and EXP034B.
  \item \textbf{A positive answer to RQ2:} aligned beats \emph{both} shuffled and
    dummy in all three experiments; the alignment-specific effect is $\Delta\in[-0.034, -0.014]$
    CE and is consistent in sign across independent training runs.
  \item \textbf{A qualified answer to RQ3:} MIR probing recovers weak key and
    pitch-class attributes well above the zero-feature baseline, while meter and
    instrument micro-F1 are excluded as prior-dominated.
  \item Two negative results (EXP027, EXP028) that establish stream-native interface
    as a necessary design constraint.
\end{enumerate}

%% ============================================================
\section{Related Work}
%% ============================================================

\subsection{Audio Language Models}

AudioLM~\citep{borsos2023audiolm} pioneered hierarchical discrete token prediction for
long-form audio continuation.  AudioPaLM~\citep{rubenstein2023audiopalm} unified speech
and audio tasks under a single LLM.  UniAudio~\citep{yang2024uniaudio} extended this to
a broad task catalog using a multi-stream tokenizer.  UniAudio 2.0~\citep{yang2026uniaudio2}
introduced ReasoningCodec, which factorizes audio into reasoning tokens (high-level,
language-aligned) and reconstruction tokens (fine-grained acoustic detail), supervised
via SFT and GRPO.  UALM~\citep{ualm2025} uses a Qwen2.5-based decoder with continuous
audio embeddings and X-codec tokens, demonstrating competitive MMAU performance but
without explicit symbolic grounding.

None of these systems incorporate a learned \emph{symbolic} planning stream verified by
alignment controls.  Our work fills this gap on the music domain.

\subsection{Symbolic Music Representations}

Symbolic music can be represented via raw MIDI events, piano rolls, or higher-level
abstractions~\citep{fradet2021miditok,hawthorne2019enabling}.  Raw MIDI tokenizations
suffer from extreme sequence lengths and are not directly comparable to audio token rates.
Compact representations such as chord sequences or beat annotations are hand-designed and
limited in expressiveness.  \splan{} learns a compact representation end-to-end from
paired MIDI features, avoiding hand-design while keeping sequence length tractable.

\subsection{Vector Quantization for Audio}

Residual vector quantization (RVQ) has become the standard for neural audio
codecs~\citep{zeghidour2021soundstream,defossez2022high,kumar2024highfidelity}.
We apply the same quantization principle to symbolic features rather than waveforms,
enabling \splan{} to use the same multi-stream position encoding as the acoustic codec.

\subsection{Music Information Retrieval}

MIR tasks including key estimation, beat tracking, instrument recognition, and chord
detection provide objective, structured evaluation targets~\citep{raffel2014mir_eval}.
We use MIR attributes from Slakh MIDI metadata as supervision and evaluation targets for
\splan{}, yielding interpretable quality measures beyond cross-entropy.

%% ============================================================
\section{Background: UniAudio 2.0}
%% ============================================================

UniAudio 2.0~\citep{yang2026uniaudio2} organizes audio into two factorized streams:

\begin{description}
  \item[Reasoning tokens $\Ramath$:] Low-rate, language-aligned tokens from a
    ReasoningCodec encoder, trained with SFT + GRPO on ASR, captioning, and
    audio-reasoning tasks.  These tokens capture high-level semantic and contextual
    content.
  \item[Reconstruction tokens $\Camath$:] An 8-codebook RVQ stream for fine-grained
    acoustic reconstruction and generation, trained after the reasoning branch is frozen.
\end{description}

The unified autoregressive model predicts $\Camath$ conditioned on text and $\Ramath$,
using a single transformer with multi-stream positional encoding.
Audio-specific experts update only audio-position hidden states.
This factorization allows both understanding (via $\Ramath$) and generation
(via $\Camath$) within a single framework.

Our work asks: \emph{can a symbolic planning stream $\splanmath$, inserted between
$\Ramath$ and $\Camath$, carry music-structural information that $\Ramath$ alone does
not provide?}

%% ============================================================
\section{The \texorpdfstring{\splan{}}{S-plan} System}
%% ============================================================

\subsection{MIDI Feature Extraction}

Given a 30-second Slakh audio-MIDI window, we extract frame-level features from the
aligned MIDI file at 100 ms resolution:

\begin{itemize}
  \item \textbf{Binary features:} Instrument-class presence vector (34 Slakh classes),
    12-dim pitch-class binary, beat/downbeat phase indicator.
  \item \textbf{Continuous features:} Instantaneous tempo (BPM), note density
    (events/second), polyphony count, instrument-weighted energy proxy.
\end{itemize}

Features are normalized per-dimension across the training split.  The final frame
vector has 53 dimensions (41 binary, 12 continuous).

\subsection{Symbolic RVQ Codec}

We train a residual vector quantizer over the frame-level feature sequence:
\begin{itemize}
  \item \textbf{Encoder:} Two-layer 1D CNN followed by temporal pooling that maps
    variable-length feature sequences to a fixed 30-frame representation per
    30-second window.
  \item \textbf{Codebook:} 4 RVQ stages, 512 codes each, trained with exponential
    moving average updates and codebook restart.
  \item \textbf{Decoder:} Symmetric 1D CNN that reconstructs the original feature
    sequence from quantized codes.
  \item \textbf{Loss:} Binary cross-entropy for binary feature dimensions plus mean
    squared error (MSE) for continuous dimensions.
\end{itemize}

Training runs for 42{,}000 steps on the Slakh training split.

\subsection{Stream-Native Integration}

We integrate \splan{} into the UniAudio 2.0 sequence as a third stream:
\[
\underbrace{[\mathrm{text}]}_{\text{1 stream}} \;\to\;
\underbrace{[\Ramath]}_{\text{8 streams}} \;\to\;
\underbrace{[\splanmath]}_{\text{4 streams}} \;\to\;
\underbrace{[\Camath]}_{\text{8 streams}}
\]
The layout produces 151 reason frames, 30 symbolic frames, and 376 semantic frames
(total 557 positions), fitting within the 1024 context limit with 0\% truncation.

We use LoRA adapters~\citep{hu2022lora} on the UniAudio 2.0 transformer attention
modules for parameter-efficient fine-tuning (140 modules, ${\sim}10.8$M trainable
parameters in the reason-only configuration and ${\sim}253$M for symbolic stream
variants), with a learned scalar gate that controls symbolic stream influence magnitude.

\subsection{Why Not Raw MIDI?}

Our design is motivated by three failures of raw MIDI appending:
\begin{itemize}
  \item \textbf{Length:} Standard MIDITok~\citep{fradet2021miditok} tokenizations
    of 30-second windows average ${\sim}2{,}000$--$8{,}000$ tokens, saturating the
    1024 context.
  \item \textbf{Alignment:} Raw MIDI event timestamps are asynchronous with audio
    frames, creating positional artifacts incompatible with multi-stream encoding.
  \item \textbf{Controllability:} Variable-length event sequences do not permit
    fixed-length shuffled/dummy controls, making the evaluation design impossible.
\end{itemize}

Low-rate rule summaries (hand-crafted text-like symbolic summaries, ${\sim}149$ tokens)
avoid the length problem but remain hand-designed.  \S\ref{sec:rq1} shows that learned
\splan{} outperforms rule summaries in acoustic-token utility.

%% ============================================================
\section{Experimental Setup}
%% ============================================================

\subsection{Dataset}
\label{sec:dataset}

All experiments use Slakh2100~\citep{manilow2019slakh}, a multi-track paired
audio-MIDI dataset synthesized with professional-grade sample-based virtual instruments.
We split at the \emph{track} level before windowing to prevent train-test contamination:

\begin{table}[h]
\centering
\small
\begin{tabular}{lrrr}
\toprule
Split & Tracks & 30 s windows & Notes \\
\midrule
Train      & 1,289 & 11,279 & Track-level split \\
Validation & 270   & 2,334  & Zero track overlap \\
Test       & 151   & 1,385  & Held-out tracks \\
\bottomrule
\end{tabular}
\caption{Slakh dataset statistics.  All symbolic comparisons use the same split.}
\label{tab:dataset}
\end{table}

\subsection{Control Protocol}

For every experiment that uses \splan{}, we train and evaluate three symbolic
conditions under matched hyperparameters:
\begin{description}
  \item[Aligned \splan{}:] Tokens from the audio window's own MIDI file, preserving
    temporal and content alignment.
  \item[Shuffled \splan{}:] Tokens of a randomly selected \emph{different} window,
    preserving token statistics but breaking content alignment.
  \item[Dummy stream:] Zero or uniformly random RVQ codes, providing a token-budget
    control without any symbolic content.
\end{description}

A gain of aligned over \emph{both} shuffled and dummy is required to attribute the
improvement to alignment-specific content (RQ2).

\subsection{Metrics}

\begin{itemize}
  \item \textbf{Cross-entropy (CE) / perplexity (PPL):} Token-level prediction loss
    on $\Camath$ acoustic tokens.
  \item \textbf{Token accuracy:} Exact-match accuracy on held-out $\Camath$ tokens.
  \item \textbf{Symbolic gate:} The learned gate value, measuring the fraction of
    symbolic stream influence admitted to the acoustic prediction path.
  \item \textbf{MIR metrics (EXP031A):} Instrument macro/micro-F1, pitch-class
    macro-F1, meter accuracy, weak key accuracy, segment density accuracy, and
    segment polyphony accuracy.
\end{itemize}

%% ============================================================
\section{Results}
%% ============================================================

We present results organized around the three research questions.
\S\ref{sec:codec} first establishes that the symbolic codec is technically sound ---
a necessary precondition for the RQ1--RQ3 claims.
\S\ref{sec:rq1} and \S\ref{sec:rq2} jointly address RQ1 (acoustic utility) and
RQ2 (alignment specificity) across four experiments.
\S\ref{sec:rq3} addresses RQ3 (music-structural content) via MIR probing.
\S\ref{sec:interface} documents the two negative results that establish stream-native
interface as a design constraint.

\subsection{Prerequisite: Codec Reconstruction Quality}
\label{sec:codec}

Before asking whether \splan{} carries useful symbolic information downstream, we verify
that the RVQ codec faithfully reconstructs the input symbolic features.

\begin{table}[h]
\centering
\small
\begin{tabular}{lc}
\toprule
Metric & Value \\
\midrule
Validation reconstruction loss & 0.00918 \\
Binary feature accuracy & 0.99777 \\
Continuous feature MSE & 0.00172 \\
Codebook utilization (books 1--4) & 31.3\% / 71.9\% / 80.7\% / 81.2\% \\
\bottomrule
\end{tabular}
\caption{Codec reconstruction quality at step 42,000.
  Per-window binary accuracy is $1.000$ on all three validation fixed examples.
  Books 2--4 reach 72--81\% utilization: no codebook collapse.}
\label{tab:exp024}
\end{table}

The codec produces a technically healthy symbolic bottleneck: it compresses
MIDI-derived structure without collapse and reconstructs features faithfully.
This is a prerequisite, not a claim --- it establishes that differences in downstream
experiments are attributable to symbolic content, not codec failure.

\subsection{RQ1: Aligned \splan{} Reduces Acoustic Token Uncertainty}
\label{sec:rq1}

We present three experiments in order of increasing rigor, each demonstrating
that aligned \splan{} reduces $\Camath$ cross-entropy below reason-only.

\paragraph{EXP025 --- Proxy prediction task (first positive signal).}
EXP025 trains a small transformer head to predict $\Camath$ tokens from frozen $\Ramath$ +
symbolic representations.  This is a proxy task using the frozen codec output rather than
end-to-end integration.

\begin{table}[h]
\centering
\small
\begin{tabular}{lccc}
\toprule
Condition & Valid CE & Test CE & Test PPL \\
\midrule
Reason only            & 7.5650 & 7.6573 & 2116.1 \\
Low-rate rule summary  & 7.5534 & 7.6445 & 2089.2 \\
Aligned \splan{}       & \textbf{7.5073} & \textbf{7.5963} & \textbf{1990.9} \\
\quad$\hookrightarrow$ Shuffled \splan{}  & 7.5376 & 7.6307 & 2060.6 \\
\quad$\hookrightarrow$ Dummy stream       & 7.6023 & 7.6942 & 2195.5 \\
\quad$\hookrightarrow$ Random stream      & 7.5400 & 7.6325 & 2064.2 \\
\bottomrule
\end{tabular}
\caption{EXP025: proxy task, 5,000 steps, 1,385 held-out test windows.}
\label{tab:exp025}
\end{table}

Aligned \splan{} achieves the lowest test CE ($7.5963$), outperforming reason-only
($-0.061$), low-rate rule summaries ($-0.048$), shuffled ($-0.034$), and dummy ($-0.098$).
This is the first evidence that \splan{} carries acoustically useful information beyond
what $\Ramath$ already encodes.

\paragraph{EXP029 --- Stream-native warmup (replication with correct interface).}
EXP029 implements the full three-stream layout with streams aligned to the base model's
positional encoding (3,000 steps, seed 260630).

\begin{table}[h]
\centering
\small
\begin{tabular}{lcccc}
\toprule
Condition & Test CE & Test Acc & Prefix-8 CE & Gate \\
\midrule
Aligned \splan{}                          & \textbf{5.766} & \textbf{0.1097} & \textbf{5.982} & 0.0033 \\
\quad$\hookrightarrow$ Shuffled \splan{} & 5.785 & 0.1067 & 6.000 & 0.0032 \\
\quad$\hookrightarrow$ Dummy stream      & 5.785 & 0.1068 & 6.004 & 0.0034 \\
\bottomrule
\end{tabular}
\caption{EXP029: stream-native layout, 3,000 steps.  Sequence length 557 tokens, 0\% truncation.}
\label{tab:exp029}
\end{table}

Aligned achieves test CE $5.766$ vs.\ shuffled $5.785$ ($\Delta = -0.018$) and dummy
$5.785$ ($\Delta = -0.019$).  The margin is modest but consistent.

\paragraph{EXP034B --- Hardened four-condition comparison (primary result).}
EXP034B is the most rigorous test of RQ1.  All four conditions run under identical
hyperparameters (800 steps, batch 1, gradient accumulation 16, LoRA lr $10^{-5}$,
symbolic lr $3\times10^{-4}$, gate lr $2\times10^{-3}$, seed 260635), differing only
in the symbolic stream variant.

\begin{table}[h]
\centering
\small
\begin{tabular}{lcccc}
\toprule
Condition & Test CE & Test PPL & Test Acc & Gate \\
\midrule
Reason-only baseline               & 5.859 & 350.4 & 0.1039 & 0.000 \\
Aligned \splan{} stream            & \textbf{5.600} & \textbf{270.4} & \textbf{0.1207} & 0.177 \\
\quad$\hookrightarrow$ Shuffled \splan{} & 5.614 & 274.3 & 0.1191 & 0.167 \\
\quad$\hookrightarrow$ Dummy stream      & 5.614 & 274.3 & 0.1193 & 0.163 \\
\bottomrule
\end{tabular}
\caption{EXP034B: all four conditions complete at step 800 (\texttt{claim\_ready = true}).}
\label{tab:exp034b}
\end{table}

Aligned achieves test CE $5.600$, substantially below reason-only ($-0.259$) and below
both corrupted controls ($-0.014$).  This answers \textbf{RQ1 positively}: aligned
\splan{} consistently reduces acoustic token uncertainty across three independent experiments.

\subsection{RQ2: The Benefit Is Alignment-Specific, Not Token Capacity}
\label{sec:rq2}

Any fixed-length stream added to $\Ramath$ expands the context available for $\Camath$
prediction.  RQ2 asks whether the gain in \S\ref{sec:rq1} is attributable to this
capacity effect, or to the aligned content specifically.

\paragraph{Decomposing the EXP034B gain.}
The $\Delta = -0.259$ CE improvement of aligned over reason-only decomposes as:

\begin{itemize}
  \item \textbf{Stream capacity effect:} Any symbolic stream variant outperforms
    reason-only by $\Delta \approx -0.245$ CE (PPL: $350 \to 274$), reflecting
    the additional context length regardless of content.
  \item \textbf{Alignment-specific effect:} Aligned \splan{} outperforms both
    shuffled ($\Delta = -0.014$) and dummy ($\Delta = -0.015$).  The shuffled
    condition preserves token statistics while breaking temporal alignment; the
    dummy condition removes all symbolic content.  Aligned beating \emph{both}
    isolates content alignment as the source of the additional gain.
\end{itemize}

\paragraph{Consistency across experiments.}

\begin{table}[h]
\centering
\small
\begin{tabular}{lccc}
\toprule
Experiment & Training steps & Setup & $\Delta$ aligned vs.\ shuffled \\
\midrule
EXP025  & 5,000 & Proxy prediction head & $-0.034$ \\
EXP029  & 3,000 & Stream-native warmup  & $-0.018$ \\
EXP034B & 800   & Full hardening        & $-0.014$ \\
\bottomrule
\end{tabular}
\caption{Alignment-specific gain across three independent experiments.  Consistent sign
  across training regimes rules out a single-run artifact.}
\label{tab:rq2_summary}
\end{table}

The alignment-specific margin decreases from EXP025 to EXP034B, likely because EXP034B
uses a shorter training budget where all conditions are closer to initialization.
The direction is consistent, providing convergent evidence across independent experimental
setups.

\paragraph{Gate behavior as a corroborating diagnostic.}
In EXP034B, the learned gate is higher for aligned ($0.177$) than shuffled ($0.167$)
or dummy ($0.163$).  The model learns to admit more information from the aligned stream,
consistent with it carrying useful content.  This corroborates the CE evidence through
an independent signal.

Together, these results answer \textbf{RQ2 positively}: the alignment-specific effect
is small but consistent and corroborated by the gate.  It is not explainable by
token capacity alone.

\subsection{RQ3: \splan{} Encodes Music-Structural Attributes}
\label{sec:rq3}

RQ3 asks whether \splan{} embeddings contain music-structural information recoverable
by supervised probing.  We train linear probing heads on frozen \splan{} features and
compare against a zero-feature baseline that reveals which axes are prior-dominated
(i.e., explained by head bias and BatchNorm statistics without any feature input).

\begin{table}[h]
\centering
\small
\begin{tabular}{lcc}
\toprule
MIR attribute & Aligned \splan{} & Zero-feature \\
 & (test / step 625) & (valid / step 500) \\
\midrule
Instrument macro-F1  & 0.501 & 0.426 \\
Instrument micro-F1  & 0.791 & \textbf{0.835} \\
Pitch-class macro-F1 & \textbf{0.854} & 0.830 \\
Meter accuracy       & 0.783 & \textbf{0.855} \\
Weak key accuracy    & \textbf{0.557} & 0.126 \\
Segment density acc. & 0.874 & \textbf{0.879} \\
Segment polyphony acc. & \textbf{0.761} & -- \\
\bottomrule
\end{tabular}
\caption{EXP031A MIR probing.  Bold marks the row where the feature-dependent model
  is more robustly attributable.  Label-shuffle and target-prior controls pending.}
\label{tab:exp031a}
\end{table}

\paragraph{Axes that support a positive claim.}
Weak key accuracy ($0.557$ vs.\ $0.126$, $\Delta = +0.431$) is the strongest signal:
the zero-feature baseline is near chance, confirming that the gain is feature-specific.
Pitch-class macro-F1 ($0.854$ vs.\ $0.830$, $\Delta = +0.024$) also shows positive
feature-specific signal.  Instrument macro-F1 ($0.501$ vs.\ $0.426$, $\Delta = +0.075$)
shows some signal but requires class-imbalance controls before strong claims.

\paragraph{Axes that are prior-dominated and excluded.}
Instrument micro-F1 ($0.791$ vs.\ $0.835$, zero-feature \emph{higher}) and meter
accuracy ($0.783$ vs.\ $0.855$, zero-feature \emph{higher}) are largely explained by
common-class and common-meter priors.  These axes must be excluded from headline claims
until label-shuffle controls confirm feature-specific learning.

\textbf{RQ3 receives a qualified positive answer}: \splan{} encodes key and pitch-class
structure well above the prior baseline, while meter and instrument micro-F1 require
additional controls.

\subsection{Design Constraint: Stream-Native Interface Is Necessary}
\label{sec:interface}

The positive results in \S\ref{sec:rq1}--\S\ref{sec:rq3} depend on the stream-native
interface design.  Two earlier experiments demonstrate that naive alternatives fail.

\textbf{EXP027 (side-channel insertion):} Inserting \splan{} tokens as additional
context rows into the released UniAudio 2.0 interface produced nearly identical results
across aligned ($5.511$), shuffled ($5.512$), and dummy ($5.510$) conditions, all worse
than reason-only ($5.360$).  The model ignored symbolic content when it was injected as
an external side-channel incompatible with multi-stream positional encoding.

\textbf{EXP028 (LoRA adapter with learned gate):} Adding LoRA adapters stabilized
training (valid CE $5.412$) but the gate converged near zero ($3.36 \times 10^{-4}$),
indicating the symbolic stream was admitted at negligible magnitude.  Training stability
is not fusion evidence.

Both negative results isolate the same root cause: symbolic tokens must enter through
the same stream-native positional encoding as $\Ramath$ and $\Camath$.  These failures
are not incidental --- they motivated the design of EXP029 and EXP034B and are
reproducible under controlled conditions.

%% ============================================================
\section{Discussion}
%% ============================================================

\paragraph{What we can and cannot claim.}
The three RQs receive positive or qualified-positive answers under controlled Slakh
experiments.  The claim boundary is important: (1) all experiments use oracle MIDI
input, not predicted symbolic tokens; (2) all CE results are on semantic token
prediction, not decoded waveforms; (3) the scope is limited to Slakh synthesized audio.
These boundaries are intentional, not evasions --- they define what the current evidence
supports before the next experimental phase.

\paragraph{Capacity versus alignment specificity.}
The dominant effect in EXP034B is the stream capacity gain ($\Delta \approx -0.245$
CE), and any symbolic stream achieves this gain.  A paper reporting only the comparison
to reason-only would overstate the alignment claim.  Our multi-condition design separates
the two effects, and the alignment-specific effect ($\Delta \approx -0.014$--$-0.034$
across experiments) is the meaningful claim.  The gate corroborates it through an
independent channel.

\paragraph{Why consistency across experiments matters for RQ2.}
The alignment-specific gap is positive across EXP025 (proxy task), EXP029 (stream
warmup), and EXP034B (hardened full comparison), all with independent training runs and
different training budgets.  A finding that requires a specific seed or a specific
training length to appear is weak evidence; a finding that appears consistently in sign
across independent runs is more robust.  The decreasing magnitude is expected: EXP034B
uses the shortest training budget, so all conditions have less room to diverge.

\paragraph{What the MIR probing result adds.}
RQ3 connects the acoustic utility (RQ1/RQ2) to an interpretable mechanism.  The
strong weak-key gap ($+0.431$) suggests that \splan{} captures tonal identity as its
most robust representational axis.  The prior-dominated meter result is an honest
negative: it tells us that 30 frames/window at current training scale is insufficient to
learn meter beyond the common-time prior, and that longer sequence or stronger rhythm
supervision may be needed.

%% ============================================================
\section{Limitations}
%% ============================================================

\paragraph{No decoded generation evidence.}
All experiments evaluate cross-entropy on $\Camath$ token prediction under oracle symbolic
input.  We do not yet have waveform-level evidence that \splan{} conditioning changes
the perceptual quality or structural fidelity of decoded audio.  Generating audio from
EXP034B checkpoints and comparing aligned versus shuffled outputs perceptually remains
future work.

\paragraph{Oracle symbolic input.}
In all stream experiments, \splan{} tokens are derived from the ground-truth MIDI file.
A practical system requires either audio-to-MIDI transcription or a separate
audio-to-\splan{} predictor at inference time, neither of which is trained yet.

\paragraph{No LLM-supervised tokenizer.}
The current \splan{} is trained via reconstruction and MIR objectives but not via an
LLM decoder objective.  A language-model-supervised tokenizer (\sllm{}) would allow
text-language grounding, enabling music captioning and symbolic QA.  We treat \sllm{}
as planned future work.

\paragraph{Slakh-only scope.}
All experiments use Slakh, a synthesized paired audio-MIDI corpus.  Generalization to
real-world recordings requires either transcription or domain adaptation, outside the
current scope.

\paragraph{Pending MIR controls.}
EXP031A label-shuffle and target-prior controls are incomplete.  Instrument macro-F1
and several other MIR axes cannot be claimed as robust symbolic evidence until these
controls pass.

%% ============================================================
\section{Conclusion}
%% ============================================================

We introduced \splan{}, a compact learned symbolic planning token that compresses
MIDI-derived musical structure into a fixed-length RVQ stream inserted between the
reasoning and acoustic token streams of UniAudio 2.0.  We organized our evaluation
around three research questions and found:

\begin{itemize}
  \item \textbf{RQ1 (Acoustic Utility):} Aligned \splan{} reduces acoustic token
    cross-entropy in three independent experiments (EXP025, EXP029, EXP034B), with
    the most rigorous result being EXP034B: test CE $5.600$ vs.\ reason-only $5.859$.
  \item \textbf{RQ2 (Alignment Specificity):} Aligned \splan{} consistently beats
    both shuffled and dummy controls.  The alignment-specific effect
    ($\Delta\in[-0.034,-0.014]$ CE) is positive and consistent across all three
    experiments, corroborated by gate values that open more for aligned content.
  \item \textbf{RQ3 (Music-Structural Content):} MIR probing recovers weak key
    ($+0.431$ above zero-feature baseline) and pitch-class ($+0.024$) attributes.
    Meter and instrument micro-F1 are excluded as prior-dominated.
\end{itemize}

Two negative results (EXP027 side-channel insertion, EXP028 adapter gating) establish
stream-native interface as a necessary design constraint, not merely an engineering
preference.

The central contribution is methodological as much as empirical: we argue that
credible symbolic-fusion claims require multi-condition evaluation with matched
corrupted controls (shuffled and dummy), zero-feature MIR baselines, and honest
reporting of what each experiment can and cannot support.  The same evaluation framework
applies broadly to any future work claiming symbolic grounding in audio language models.

%% ============================================================
\section*{Acknowledgments}
%% ============================================================

This work was conducted as part of the CT-NLP course at the Graduate School of
Cultural Technology, KAIST.

%% ============================================================
\bibliography{custom_symbolic_lalm}
%% ============================================================

\appendix

\section{EXP024 Fixed-Window Gallery}
\label{app:exp024}

Table~\ref{tab:exp024_gallery} shows per-window reconstruction statistics for three
validation examples from Slakh Track01501.

\begin{table}[h]
\centering
\small
\begin{tabular}{lccp{2.6cm}}
\toprule
Window & Binary Acc & Cont. MSE & Unique codes (books 1--4) \\
\midrule
w000 (0--30 s)   & 1.000 & 0.00174 & 15 / 15 / 17 / 17 \\
w001 (30--60 s)  & 1.000 & 0.00246 & 20 / 24 / 26 / 25 \\
w003 (90--120 s) & 1.000 & 0.00130 & 17 / 17 / 21 / 25 \\
\bottomrule
\end{tabular}
\caption{EXP024 fixed validation window statistics.  All windows achieve perfect binary
  feature reconstruction.}
\label{tab:exp024_gallery}
\end{table}

\section{EXP029 Fixed-Window Case Spectrum}
\label{app:exp029}

Table~\ref{tab:exp029_cases} includes both positive and boundary cases to prevent
cherry-picking the qualitative evidence.

\begin{table}[h]
\centering
\small
\begin{tabular}{lcccc}
\toprule
Window & Aligned & Shuffled & Dummy & Margin \\
\midrule
w000 & 5.226 & 5.301 & 5.280 & $+$0.054 \\
w001 & 6.159 & 6.179 & 6.163 & $+$0.004 \\
w002 & 6.097 & 6.104 & 6.091 & $-$0.006 \\
w003 & 5.876 & 5.893 & 5.869 & $-$0.007 \\
\bottomrule
\end{tabular}
\caption{EXP029 per-window validation losses.  Margin = aligned minus best control;
  negative values are boundary cases where aligned is not the best.}
\label{tab:exp029_cases}
\end{table}

\section{EXP034B Per-Codebook Cross-Entropy}
\label{app:exp034b_cb}

Table~\ref{tab:exp034b_cb} shows per-codebook test CE for EXP034B.  The symbolic benefit
is distributed across codebooks, with the strongest gains in books 0, 2, and 3.

\begin{table}[h]
\centering
\small
\setlength{\tabcolsep}{3pt}
\begin{tabular}{lcccccccc}
\toprule
Condition & cb0 & cb1 & cb2 & cb3 & cb4 & cb5 & cb6 & cb7 \\
\midrule
Reason-only & 5.60 & 6.00 & 3.45 & 5.28 & 6.02 & 6.48 & 6.87 & 7.16 \\
Aligned \splan{} & \textbf{5.24} & 5.71 & \textbf{3.19} & \textbf{4.98} & 5.75 & 6.25 & 6.68 & 7.01 \\
Shuffled & 5.26 & 5.73 & 3.20 & 4.99 & 5.76 & 6.26 & 6.69 & 7.02 \\
Dummy & 5.26 & 5.73 & 3.20 & 4.99 & 5.76 & 6.26 & 6.69 & 7.02 \\
\bottomrule
\end{tabular}
\caption{EXP034B per-codebook test CE.  Aligned \splan{} achieves the lowest CE in
  books 0, 2, and 3.}
\label{tab:exp034b_cb}
\end{table}

\end{document}
